% ===== §6 Co-evolution as Relational Dynamics =====
\section{Co-evolution as Relational Dynamics}
\label{sec:coevolution}

\noindent
The preceding section established that a contingent event, when it enters the relational field through the constitutive act of sharing, becomes a relational event---an event whose significance belongs to the \emph{between} rather than to either individual. But what happens next? What does the relational field \emph{do} with the event it has received? The answer that the GRB framework proposes is: the relational field undergoes co-evolution. The two dynamical systems that compose the relational field---the two individual subjects, each with their own history, attunement, and dynamical structure---are perturbed by the event and respond to that perturbation together, in a way that produces changes in each that are irreducible to what either would have undergone alone.

This section develops the concept of relational co-evolution in three registers: philosophical, formal-physical, and neuroscientific. The three registers are not alternative descriptions of the same phenomenon but complementary levels of analysis---in the sense of Marr's levels \citep{marr1982}---each of which captures something that the others do not, and none of which is reducible to the others. Together, they constitute a multi-level account of how the relational field generates happiness through co-evolutionary dynamics.

\subsection{The Philosophical Register: No Isolated Individual Dynamics}
\label{subsec:philocoevo}

\noindent
The philosophical core of the co-evolutionary account is a negative claim that must be stated with precision: \emb{there are no fully isolated individual dynamics in an intimate relational system}. This is not the trivial claim that individuals influence each other---that is compatible with a picture on which each individual has their own dynamics and these dynamics are modified by external influences, including the other person. The claim is stronger: that the dynamics of each individual, within an intimate relational system, are \emph{constitutively} shaped by the relational field, so that there is no baseline individual dynamic that exists independently of the relation and to which the individual would revert if the relation were removed.

This claim follows from the ontological framework established in \xref{sec:turn}. If the individual subject is constituted by and within the relational field, then the individual's dynamical structure---their patterns of emotional response, their attentional habits, their modes of making sense of the world---is itself a product of the relational field. The self that walks alongside the beloved, notices the flower, and feels the happiness of sharing it, is not a self that exists independently of this relation and happens to be walking alongside someone; it is a self whose very capacity for this kind of noticing and sharing has been shaped by the history of this relation. The individual dynamics are, in this sense, always already relational dynamics.

The positive claim follows from the negative: when a contingent event enters the relational field and the two individuals respond to it together, the response of each is not simply their individual response to an external stimulus. It is a response that is already shaped by the relational field---by the history of shared responses, the accumulated attunement, the joint habits of attention---and that further shapes the relational field through the new history it creates. The co-evolution triggered by the shared flower is not two individual evolutions happening in parallel; it is a single relational evolution that produces changes in each individual that are functions of the relational field as a whole, not of either individual's dynamics in isolation.

\begin{claim}[The irreducibility of relational co-evolution]
The co-evolution triggered in a relational system by a contingent event is not the sum of two individual evolutions occurring in parallel. It is a single dynamical process at the level of the relational field, whose effects on each individual are functions of the field as a whole. There are no fully isolated individual dynamics in an intimate relational system: the individual dynamics are always already relational, and the co-evolutionary response to a contingent event is irreducible to any decomposition into individual components.
\end{claim}

The happiness that arises from this co-evolutionary process is, on the philosophical account, the \emph{felt signal} of the relational field's generative activity. It is the way in which the relational field's co-evolution presents itself to the individuals who participate in it---the phenomenal surface of a dynamical process that occurs at the level of the field. This is why the happiness of the shared flower feels different from any happiness that either person could generate alone: it is the felt quality of a process that exceeds both of them, that is happening at a level they participate in but do not control, and that produces in each of them a sense of something coming from the between, from the shared field, rather than from within themselves.

\subsection{The Formal-Physical Register: Coupled Dynamical Systems}
\label{subsec:formalcoevo}

\noindent
The philosophical account of co-evolution can be given formal content by drawing on the theory of coupled dynamical systems---a branch of mathematics and theoretical physics that studies the behaviour of systems that are connected to and influence each other \citep{strogatz2001,kuramoto1984}. The formal account does not replace the philosophical account but gives it mathematical precision and connects it to a rich body of theoretical and empirical results.

\subsubsection{Coupled Oscillators and the Emergence of New Attractors}
\label{subsubsec:coupled}

\noindent
Consider two dynamical systems, $\mathcal{S}_1$ and $\mathcal{S}_2$, each with their own state space and their own autonomous dynamics. In the absence of coupling, each system evolves according to its own equations of motion:
\begin{align}
\dot{\mathbf{x}}_1 &= \mathbf{f}_1(\mathbf{x}_1) \\
\dot{\mathbf{x}}_2 &= \mathbf{f}_2(\mathbf{x}_2)
\end{align}
where $\mathbf{x}_i$ is the state vector of system $i$ and $\mathbf{f}_i$ is its autonomous vector field. Each system, evolving independently, has its own attractors---its own stable states or stable cycles toward which it tends over time.

When the two systems are coupled---when each influences the other---the equations of motion become:
\begin{align}
\dot{\mathbf{x}}_1 &= \mathbf{f}_1(\mathbf{x}_1) + \varepsilon \mathbf{g}_1(\mathbf{x}_1, \mathbf{x}_2) \\
\dot{\mathbf{x}}_2 &= \mathbf{f}_2(\mathbf{x}_2) + \varepsilon \mathbf{g}_2(\mathbf{x}_1, \mathbf{x}_2)
\end{align}
where $\varepsilon$ is the coupling strength and $\mathbf{g}_i$ is the coupling function that describes how system $j$ influences system $i$. The coupled system is now a single dynamical system in the product space $\mathcal{S}_1 \times \mathcal{S}_2$, with its own dynamics, its own attractors, and its own emergent behaviour.

The crucial mathematical fact is this: the attractors of the coupled system are generally \emph{not} the attractors of the individual systems. Coupling produces new attractors---new stable states or stable cycles that exist in the product space but correspond to no attractor of either individual system. These new attractors are the formal representation of what the GRB framework calls relational co-evolution: they are the stable patterns of joint behaviour that the relational system generates through its coupling, patterns that neither individual would exhibit in isolation.

The shared flower, in this formal framework, functions as an external perturbation: a sudden change in the state of one or both systems that displaces the coupled system from its current trajectory and initiates a transient response. If the coupling is strong enough---if the relational field is deep enough---this perturbation is absorbed into the relational dynamics and processed jointly, producing a trajectory through the product space that is characteristic of the coupled system rather than of either individual. The happiness of the shared flower is, formally, the phenomenal correlate of this joint trajectory: the way in which the coupled system's response to the perturbation presents itself to the individuals who constitute it.

\subsubsection{Synchronisation and the Kuramoto Model}
\label{subsubsec:kuramoto}

\noindent
A particularly important class of coupled dynamical systems is the class of coupled oscillators: systems in which each individual component has a natural frequency of oscillation, and coupling produces synchronisation---the alignment of the individual oscillations into a common rhythm \citep{kuramoto1984}. The Kuramoto model provides the canonical mathematical treatment of this phenomenon:

\begin{equation}
\dot{\theta}_i = \omega_i + \frac{K}{N} \sum_{j=1}^{N} \sin(\theta_j - \theta_i)
\end{equation}

where $\theta_i$ is the phase of oscillator $i$, $\omega_i$ is its natural frequency, $K$ is the coupling strength, and $N$ is the number of oscillators. When $K$ exceeds a critical threshold $K_c$, the system undergoes a phase transition from incoherence to synchronisation: the individual oscillators, despite their different natural frequencies, lock into a common phase and frequency.

Applied to the relational system, the Kuramoto model captures something important about the phenomenology of shared happiness. The two individuals in an intimate relation are, in general, oscillating at different frequencies: different rhythms of attention, different emotional tempos, different cycles of engagement and withdrawal. The coupling produced by the relational field---through shared attention, shared activity, shared history---tends to synchronise these rhythms, producing a common relational rhythm that is not the rhythm of either individual but the rhythm of the coupled system.

The contingent event---the flower on the path---functions, in the Kuramoto framework, as a perturbation that temporarily increases the effective coupling strength: the shared attention to the flower aligns the two individuals' attentional rhythms in a way that ordinary co-presence may not. This transient increase in coupling strength produces a transient increase in synchronisation, and it is this transient synchronisation---the moment of aligned attention, aligned affect, aligned presence---that constitutes, at the formal level, the co-evolutionary response to the contingent event. The happiness of the shared flower is, in part, the felt quality of this transient synchronisation: the sense of being, for a moment, in rhythm with the other in a way that feels effortless and complete.

\subsubsection{Quantum Entanglement Structure and Non-Locality}
\label{subsubsec:entanglement}

\noindent
A third formal framework illuminates a different aspect of relational co-evolution: the quantum-mechanical concept of entanglement \citep{einstein1935}. We emphasise at the outset---as this paper's commitment to intellectual honesty requires---that the following is a \emb{structural isomorphism} rather than a literal physical claim. We are not asserting that intimate relations involve quantum-mechanical processes in any literal sense; we are asserting that the mathematical structure of quantum entanglement provides a precise and illuminating analogy for certain features of relational co-evolution that classical dynamical systems theory does not fully capture.

A quantum system consisting of two subsystems is said to be entangled if its state cannot be written as a product of the states of the individual subsystems:
\begin{equation}
|\Psi\rangle \neq |\psi_1\rangle \otimes |\psi_2\rangle
\end{equation}
In an entangled state, the two subsystems are correlated in a way that is not reducible to the properties of either subsystem considered independently: a measurement on one subsystem immediately affects the probabilities of measurement outcomes on the other, regardless of the spatial separation between them. The correlations are, in this sense, non-local: they are properties of the joint state rather than of either subsystem.

The structural analogy with relational co-evolution is precise. The relational field, on the GRB account, is a joint state that cannot be decomposed into the product of two individual states: the properties of the relational field are not reducible to the properties of either individual, and changes in one individual's state are immediately reflected in the state of the field as a whole, not because of any causal influence transmitted between individuals but because the individuals are, in their relational existence, parts of a joint state. The non-locality of quantum entanglement is the formal analogue of the \emph{between}'s non-reducibility to either of the individuals who constitute it.

This analogy is philosophically significant because it captures the sense in which the happiness of the shared flower is genuinely non-local: it does not reside in either person but in the relational field that both participate in, and it is experienced by each not as a property of themselves but as a property of the field in which they are jointly constituted. The happiness comes from the between; it is felt \emph{as coming from the between}; and this phenomenological fact corresponds, at the formal level, to the non-local character of the joint relational state.

We note, for the sake of the honest account that this paper's methodology requires, that the entanglement analogy faces the same limitation that \pinkref{Paper~XIII} noted in its application of quantum dynamical systems to the analysis of trust \citep{paperxiii}: the analogy illuminates the structure of the phenomenon but does not license any claims about the underlying physical implementation. The relational field is not a quantum system; its non-locality is not the non-locality of quantum mechanics. It is a structural non-locality---an irreducibility of the field to its parts---that the mathematical language of entanglement expresses with precision that ordinary language does not.

\subsection{The Neuroscientific Register: Inter-Brain Synchrony and Embodied Resonance}
\label{subsec:neuro}

\noindent
The philosophical and formal accounts of relational co-evolution are complemented, at the empirical level, by a growing body of neuroscientific and biological research on the neural and physiological correlates of shared experience. This research does not reduce the philosophical phenomenon to its neural correlates---the relationship between the formal and the physical, and between the physical and the phenomenal, remains one of the deepest unsolved problems in philosophy of mind---but it provides empirical grounding for the claim that relational co-evolution is a real, measurable phenomenon that occurs at multiple levels of the organism.

\subsubsection{Inter-Brain Synchrony: EEG, MEG, and Hyperscanning}
\label{subsubsec:hyperscanning}

\noindent
The most direct neural evidence for relational co-evolution comes from hyperscanning research: studies in which the neural activity of two or more individuals is measured simultaneously, allowing the detection of inter-brain synchrony---the alignment of neural oscillations across brains \citep{hasson2012,dikker2017}.

Electroencephalography (EEG) hyperscanning provides the highest temporal resolution: it can detect synchrony at the millisecond timescale, capturing the rapid dynamics of shared attention, joint action, and co-present affect that are most directly relevant to the phenomenology of the shared flower. EEG studies of dyadic interaction have documented inter-brain synchrony in multiple frequency bands---including theta (4--8 Hz), alpha (8--12 Hz), and gamma (30--80 Hz)---during joint attention, cooperative task performance, and face-to-face communication \citep{liu2023}. The synchrony is not merely a consequence of both brains processing the same external stimulus; it reflects genuine coupling between the two neural systems, as evidenced by the fact that it exceeds the synchrony observed when two individuals are presented with the same stimulus independently.

Magnetoencephalography (MEG) offers complementary advantages: superior spatial resolution compared to EEG (due to the absence of volume conduction distortion), and sensitivity to deep cortical sources that EEG cannot easily detect. MEG hyperscanning is technically more demanding than EEG hyperscanning---it requires magnetically shielded rooms and specialised infrastructure---but it provides a more complete picture of the spatial distribution of inter-brain synchrony, allowing the identification of the specific cortical networks involved in relational co-evolution.

Functional magnetic resonance imaging (fMRI) contributes the highest spatial resolution: it can localise the neural correlates of shared experience to specific brain regions with millimetre precision, identifying the prefrontal, temporal, and limbic structures that are differentially engaged during shared versus individual experience \citep{hasson2012}. While fMRI's temporal resolution is too low to capture the rapid dynamics of moment-to-moment co-evolution, its spatial resolution makes it indispensable for understanding the neural architecture of the relational field.

Functional near-infrared spectroscopy (fNIRS) offers a crucial practical advantage: it can be used in naturalistic environments, without the rigid constraints of MRI or MEG scanners, allowing the measurement of neural co-evolution during real-world relational interactions. fNIRS hyperscanning studies have documented inter-brain synchrony during cooperative tasks, joint music-making, and face-to-face conversation in settings that much more closely approximate the natural conditions of intimate relational life \citep{hasson2012}.

Together, these neuroimaging modalities provide convergent evidence for the neural reality of relational co-evolution. When two people share a contingent event---when they attend jointly to the flower on the path, or sit together in the wordless co-presence of relational flow---their brains are not processing the experience independently; they are coupled systems whose neural dynamics are synchronised in ways that reflect and constitute the shared experience. The inter-brain synchrony is the neural correlate of the \emph{between}: the physical instantiation, at the level of neural dynamics, of the relational field's co-evolutionary activity.

\subsubsection{Embodied Resonance and the Mirror System}
\label{subsubsec:mirror}

\noindent
Inter-brain synchrony is not confined to the cortical level; it has embodied correlates that extend throughout the organism. The discovery of mirror neurons---neurons that fire both when an organism performs an action and when it observes the same action performed by another---has provided a neural mechanism for what Merleau-Ponty described as intercorporeality: the pre-personal, pre-cognitive resonance between embodied subjects that constitutes the ground of shared experience \citep{mukamel2010}.

The mirror system is not merely a mechanism for imitation or action understanding; it is the neural substrate of embodied resonance---the way in which the other's bodily states, expressions, and actions are directly registered in one's own body, not as representations of the other's states but as activations of one's own motor and emotional systems. When the beloved is moved by the flower, the lover's mirror system registers this movement not as an observation of an external event but as a resonance within the lover's own embodied being. The shared emotion is not communicated from one person to another; it arises, simultaneously and through mutual embodied resonance, in both.

This embodied resonance is the neural and phenomenological correlate of what the formal framework describes as coupling: the mutual influence of two dynamical systems that produces co-evolutionary dynamics in the product space. The coupling function $\mathbf{g}_i(\mathbf{x}_1, \mathbf{x}_2)$ in the formal account is, at the neural level, implemented by the mirror system and the broader embodied resonance mechanisms; and the co-evolutionary trajectory through the product space is, at the phenomenological level, the shared emotional experience that arises from this mutual embodied activation.

\subsubsection{Co-regulation and the Neurobiology of Attachment}
\label{subsubsec:coregulation}

\noindent
A third line of neuroscientific evidence for relational co-evolution comes from the study of physiological co-regulation: the finding that intimate partners mutually regulate each other's physiological states, including heart rate, cortisol levels, respiratory rhythm, and autonomic nervous system tone \citep{feldman2007,porges2011}.

Co-regulation is most dramatically documented in the parent-infant dyad, where the caregiver's physiological state continuously influences the infant's---and, through the infant's responsive signals, is itself influenced by the infant's state \citep{feldman2007}. But co-regulation extends throughout the lifespan and is robustly documented in adult intimate relationships: partners synchronise their heart rate variability during shared emotional experiences, their cortisol levels covary in response to shared stressors, and their autonomic regulation is mutually dependent in ways that make each partner's physiological well-being genuinely dependent on the other's physiological state.

This physiological co-regulation is not a mere consequence of shared circumstances---both partners being exposed to the same environmental conditions---but a genuine coupling of two biological systems, mediated by the full range of interpersonal signals (vocal prosody, facial expression, touch, gaze, bodily posture) that constitute the embodied communication channel of an intimate relation. The two partners are, in a physiologically precise sense, coupled dynamical systems: each partner's physiological dynamics are constitutively shaped by the other's, so that the system they compose together has properties---including regulatory capacities and resilience---that neither possesses alone.

The happiness of the shared flower has, on this account, a physiological dimension that is irreducible to either partner's individual physiological state: it is the felt quality of a moment of successful co-regulation, a moment in which the two physiological systems are in a particularly harmonious alignment---a joint homeostatic state, reached through the mutual coupling of embodied dynamics, that neither could achieve alone.

\begin{claim}[Multi-level co-evolution]
Relational co-evolution is not a metaphor or a philosophical abstraction. It is a real, multi-level process that occurs simultaneously at the level of conscious experience (the phenomenology of shared happiness), neural dynamics (inter-brain synchrony in EEG, MEG, fMRI, and fNIRS), embodied resonance (the mirror system and intercorporeal activation), and physiological regulation (heart rate, cortisol, and autonomic co-regulation). These levels are not alternative descriptions but genuinely distinct levels of the same phenomenon, each capturing something the others do not, and together constituting the full structure of relational co-evolution.
\end{claim}

The multi-level character of relational co-evolution is itself philosophically significant. It means that the happiness of the shared flower is not a purely mental event, not a purely neural event, not a purely physiological event, but an event that occurs at all of these levels simultaneously and that can only be fully understood by attending to all of them. This is consistent with the GRB framework's commitment to a non-reductive, multi-register analysis of relational phenomena: the formal, the neural, the embodied, and the phenomenological are all genuine levels of the relational field's co-evolutionary activity, and the philosophy of relational happiness must be answerable to all of them.
